\documentclass[10pt]{article}

\usepackage[utf8]{inputenc}

\usepackage[T1]{fontenc}

\usepackage{amsmath}

\usepackage{amsfonts}

\usepackage{amssymb}

\usepackage[version=4]{mhchem}

\usepackage{stmaryrd}

\usepackage{mathrsfs}

\usepackage{graphicx}

\usepackage[export]{adjustbox}

\graphicspath{ {./images/} }

\usepackage{bbold}



\begin{document}

СТОКСА ФОРМУЛА - 1) формула, выражающая связь между потоком векторного поля через двумерное ориентированное многообразие и циркуляцию этого поля по соответствующим образом ориевти рованному краю этого многообразия. Пусть $S-$ ориентированная кусочно гладкая поверхность, $v=(\cos \alpha, \cos \beta, \cos \gamma)-$ единичная нормаль к поверхности $S$ (в тех точках, конечно, где она существует), задающая ориентацию $S$, и пусть край поверхности $S$ состоит из конечного числа кусочно гладких контуров. Через $\partial S$ обозначен край поверхности $S$, ориентированный с помощью единичного касательного к нему вектора $\tau$ так, чтобы получающаяся ориентация края $\partial S$ была согласована с ориентацией $v$ поверхности $S$.



Если $\boldsymbol{a}=(P, Q, R)$ - непрерывно дифференцируемое в окрестности поверхности $S$ векторное поле, то



\[

\iint_{S}(\operatorname{rot} \boldsymbol{a}, \boldsymbol{v}) d S=\int_{\partial S}(\boldsymbol{a}, \boldsymbol{\tau}) d s

\]



( $d S$ - элемент площади поверхности $S, d s$ - дифференциал длины дуги края $\partial S$ поверхности $S$ ) или, в координатном виде:



\[

\iint_{S}\left|\begin{array}{ccc}

\cos \alpha & \cos \beta & \cos \gamma \\

\frac{\partial}{\partial x} & \frac{\partial}{\partial y} & \frac{\partial}{\partial z} \\

P & Q & R

\end{array}\right| d S=\int_{\partial S} P d x+Q d y+R d z .

\]



Предложена Дж. Стоксом (G. Stokes, 1854). 2)C. ф. наз. также обобщение формулы (*), представляющее собой равенство интеграла от внешнего дифференциала дифференциальной формы $\omega$ по ориентированному компактному многообразию $M$ и интеграла от самой формы $\omega$ по ориентированному согласованно с ориентацией многообразия $M$ краюо $\partial M$ многообразия $M$ :



\[

\int_{M} d \omega=\int_{\partial M} \omega \text {. }

\]



Частными случаями этой формулы являются Ньютона - Лейбница формула, Грина формула, Остроградского формула.



СТОКСА ЯВЛЕНИЕ - свойство функции $f(z)$ иметь различные асимптотические выражения при $|z| \rightarrow \infty$ в различных областях комплексной плоскости z. Дж. Стокс показал [1], что решение $w_{0}(z)$ т. н. уравнения Эั̆ри:



\[

w^{\prime \prime}-z w=0,

\]



убывающее при действительных $z=x \rightarrow+\infty$, имеет асимптотику шри $|z| \rightarrow \infty$ :



\[

\begin{gathered}

w_{0}(z) \sim C_{z^{-1 / 4}} \exp \left(-\frac{2}{3} z^{3 / 2}\right), \\

|\arg z| \leqslant \pi-\varepsilon<\pi ; \\

w_{0}(z) \sim C e^{i \pi / 4} z^{-1 / 4} \cos \left(\frac{2}{3} z^{3 / 2}-\frac{\pi}{4}\right), \\

|\arg z-\pi| \leqslant \varepsilon<\pi,

\end{gathered}

\]



где $C \neq 0$ - постоянная; функция $w_{0}(z)$ - целая, а ее асимптотика - разрывная функция.



С. я. имеет место для интегралов ЈІапласа, решений обыкновенных диффференциальных уравнений и т. д. (cM. [2], [3])



Jum.: [1] S t o k e s G. G., "Trans. Cambridge Phil. Soc.», 1864, v. 10 , p. $106-28 ;$ [2] Х е д и н г Д ж., Введение в метод фазовых интегралов (Метод ВКБ), пер. с англ., М., 1965; [3] Б р ё й н д е Н. Г., Асимптотические методы в анализе, пер.



СТОУНА ПРОСТРАНСТВО булевой алгебры $\mathscr{B}-$ вполне несвязное бикомпактное пространство $(X, \tau)$, поле всех открыто-замкнутых множеств к-рого изоморф-



\includegraphics[max width=\textwidth, center]{2023_07_09_2a337271decbbf4822c4g-1(1)}

$\mathscr{B}$ следующим образом: $X$ есть множество всех ультрафильтров $\mathscr{A}$, а топология $\tau$ порождена семействой подмножеств вида $\{\xi \in X: A \in \xi\}$, где $A$ - произвольный



\includegraphics[max width=\textwidth, center]{2023_07_09_2a337271decbbf4822c4g-1}

вать множества максимальных идеалов, двузначных гомоморфизмов, двузначных мер на $\mathscr{B}$ с соответствующей топологией. Изоморфные булевы алгебры имеют гомеоморфные С. п. Каждое вполне несвязное бикомпактное пространство есть С. Ш. булевой алгебры своих открыто-замкнутых множеств.



Понятие С. п. и основные его свойства найдены п исследованы М. Стоуном (M. Stone, 1934-37, cм. [1]).



C. п. булевой алгебры метризуемо тогда и только тогда, когда она счетна. Булева алгебра полна тогда и только тогда, когда ее С. п. экстремально несвязно (т. е. замыкание любого открытого множества в нем открыто). Канторово совершенное множество есть С. І. счетной безатомной бесконечной булевой алгебры (все они изоморфны). Канторов обобщенный дисконтинуум $D^{m}$ есть C. П. свободной булевой алгебры с $m$ образующими. Лит.: [1] С и к о р с к и й Р., Булевы алгебры, пер. с англ., СТОУНА РЕIIIETКА - дистрибутивная решетка $L$ с псевдодополнениями (см. Решетка с дополнениями), в к-рой $a^{*}+a^{* *}=1$ для всех $a \in L$. Дистрибутивная решетка $L$ с псевдодополнениями является С. р. тогда и только тогда, когда теоретико-структурное объединение двух ее различных минимальных простых идеалов совпадает с $L$ (т е о р е м а $\Gamma$ р е т п е р а T a, [3]).



C. p., рассматриваемая как универсальная алгебра с основными операциями $\left\langle\vee, \wedge,{ }^{*}, 0,1\right\rangle$, наз. а л г е бр о й С т о ун а. Всякая алгебра Стоуна является подпрямым произведением двухэлементных и трехәлементных цепей. В решетке с псевдодополнениями элемент $x$ наз. п л о т н ы м, если $x^{*}=0$. Центр $C(L)$ репетки Стоуна $L$ - булева алгебра, а множество $D(L)$ всех ее плотных әлементов - дистрибутивная решетка с единицей. При этом гомоморфизм $\varphi^{L}$ решетки $C(L)$ в решетку $F(D(L))$ фильтров решетки $D(L)$, определяемый условием



\[

a \varphi^{L}=\left\{x \mid x \in D(L), x \geqslant a^{*}\right\},

\]



сохраняет 0 и 1.



Т р о й к о й, ассоциироваңной с алгеброй Стоуна $L$, наз. тройка $\left\langle C(L), D(L), \varphi^{L}\right\rangle$. Естественным образом определяются гомоморфизмы и изоморфизмы троек. Произвольная тройка $\langle C, D, \varphi\rangle$, где $C$ - булева алгебра, $D-$ дистрибутивная решетка с 1 , a $\varphi: C \rightarrow F$ (D) - гомоморфизм, сохраняющий 0 и 1, изоморфна тройке, ассоциированной с нек-рой алгеброй Стоуна; алгебры Стоуна изоморфны тогда и только тогда, когда изоморфны ассоциированные с ними тройки (т е о р ем а Чен а - Гр е тц е р a, [2]).



Лит.: [1] Б и р к го ф Г., Теория решеток, пер. с англ., М., 1984; [2] C h e n C. C., Gr ätz e r G., «Canad. J. Math.", E. T., "Acta math. Acad. sci. hung.», 1957, v. 8, fasc. 3-4, p. $455-60 ;[4]$ Ф о ф а н о в а Т. С., в кн.: Упорядоченные множества и решетки, в. 3 , Саратов, 1975 , с. $22-40$.



CTОУНА - ЧЕХА БИКОМПАКТНОЕ РАСII ИРЕНИЕ - наибольшее бикомпактное расширение $\beta X$ вполне регулярного топологич. пространства $X$. Построено Э. Чехом [1] и М. Стоуном [2].



Пусть $\left\{f_{\alpha}: X \rightarrow[0,1]\right\}_{\alpha \in A}-$ множество всех непрерывных функций $X \rightarrow[0,1]$. Отображение $\varphi: X \rightarrow \mathbb{R} A$, где $\varphi(x)_{\alpha}=f_{\alpha}(x)$, является гомеоморфизмом на свой образ. Тогда, по определению, $\beta X=[\varphi(X)]$ (где [ $\cdot]$ операция замыкания). Для любого бикомпактного расширения $b X$ существует непрерывное отображение $\beta X \rightarrow b X$, тождественное на $X$, что и выражается эпитетом «наибольшее».



С. - Ч. б. р. квазинормального пространства совпадает с его Уолмена бикомпактным расширением.



Jum.: [1] C e c h E., "Ann. Math.», 1937, v. 38, p. 823-44;

St on e M., "Trans. Amer. Soc.», 1937, v. 41. p. 375-





\end{document}